\section*{Hoofdstuk \ref{ch:g4:mult} --- Grotere getallen vermenigvuldigen}

\begin{solution}{exo:g4:mult:1}
$470$; \quad $2\,300$; \quad $720$; \quad $4\,500$; \quad
$2\,100$.
\end{solution}

\begin{solution}{exo:g4:mult:2}
$34 \times 21 = 34 \times 20 + 34 = 680 + 34 = 714$; het cijferen geeft
dezelfde regels: $34$ en $680$.
\end{solution}

\begin{solution}{exo:g4:mult:3}
$63 \times 24 = 1\,512$; \quad
$85 \times 47 = 3\,995$; \quad
$126 \times 53 = 6\,678$.
\end{solution}

\begin{solution}{exo:g4:mult:4}
$208 \times 34$: schatting $200 \times 34 = 6\,800$; precies $7\,072$.

$517 \times 62$: schatting $500 \times 60 = 30\,000$; precies
$32\,054$.
\end{solution}

\begin{solution}{exo:g4:mult:5}
$45 \times 11 = 450 + 45 = 495$.

$32 \times 99 = 3\,200 - 32 = 3\,168$.

$24 \times 45 = 12 \times 90 = 1\,080$ (of nog een keer:
$6 \times 180 = 1\,080$).
\end{solution}

\begin{solution}{exo:g4:mult:6}
$26 \times 18 = 468$ stoelen.
\end{solution}

\begin{solution}{exo:g4:mult:7}
Schatting: $30 \times 50 = 1\,500$. Precies: $32 \times 48 = 1\,536$
kratten.
\end{solution}

\begin{solution}{exo:g4:mult:8}
Ze is vergeten dat de $2$ van $23$ \emph{tientallen} telt: de tweede regel
moet $57 \times 20 = 1\,140$ zijn, niet $114$. Het juiste totaal is
$171 + 1\,140 = 1\,311$.
\end{solution}

\begin{solution}{exo:g4:mult:9}
Tulpen: $14 \times 35 = 490$. Narcissen: $12 \times 28 = 336$. Samen:
$490 + 336 = 826$ bloemen.
\end{solution}

\begin{solution}{exo:g4:mult:10}
Schatting: $70 \times 30 = 2\,100$ kg. Precies: $67 \times 30 = 2\,010$
kg.
\end{solution}

\begin{solution}{exo:g4:mult:11}
Een rechthoek van $15 \times 12$ vakjes kun je met een verticale lijn in
een stuk $15 \times 10$ en een stuk $15 \times 2$ knippen: elk stuk tellen
geeft $150 + 30 = 180$. Een andere knip werkt net zo goed, bijvoorbeeld
$15 = 10 + 5$: $10 \times 12 + 5 \times 12 = 120 + 60 = 180$. Hoe je de
rechthoek ook knipt, het totale aantal vakjes blijft hetzelfde.
\end{solution}
